% !TeX program = lualatex
\documentclass[12pt,aspectratio=169]{beamer}
\usepackage[utf8]{inputenc}
\usepackage[T1]{fontenc}
\usepackage{amsmath,amsfonts,amssymb}
\usepackage{graphicx}
\usepackage{tikz}
\usepackage{xcolor}
\usepackage{hyperref}
\usepackage{anyfontsize}
\usepackage{lmodern}

% ====== MAKE ALL FONTS SMALLER ======
\renewcommand{\tiny}{\fontsize{5}{6}\selectfont}
\renewcommand{\scriptsize}{\fontsize{6}{7}\selectfont}
\renewcommand{\footnotesize}{\fontsize{7}{8}\selectfont}
\renewcommand{\small}{\fontsize{8}{9}\selectfont}
\renewcommand{\normalsize}{\fontsize{9}{10}\selectfont}
\renewcommand{\large}{\fontsize{10}{11}\selectfont}
\renewcommand{\Large}{\fontsize{11}{13}\selectfont}
\renewcommand{\LARGE}{\fontsize{13}{15}\selectfont}
\renewcommand{\huge}{\fontsize{15}{17}\selectfont}
\renewcommand{\Huge}{\fontsize{18}{21}\selectfont}

% ====== COLOR THEME ======
\definecolor{redcorrect}{RGB}{200,30,30}
\definecolor{bluecorrect}{RGB}{30,80,200}
\definecolor{greencheck}{RGB}{0,150,50}
\definecolor{lightgray}{RGB}{245,245,245}
\definecolor{darkgray}{RGB}{50,50,50}
\definecolor{softyellow}{RGB}{255,248,200}
\definecolor{infoblue}{RGB}{230,240,252}
\definecolor{steppurple}{RGB}{150,50,200}
\definecolor{deepblue}{RGB}{20,60,150}

\setbeamercolor{title}{fg=white,bg=redcorrect}
\setbeamercolor{frametitle}{fg=white,bg=redcorrect}
\setbeamercolor{block title}{fg=white,bg=bluecorrect}
\setbeamercolor{block body}{fg=darkgray,bg=lightgray}
\setbeamercolor{normal text}{fg=darkgray}
\usecolortheme{default}

% ====== CUSTOM COMMANDS ======
\newcommand{\correct}[1]{\textcolor{greencheck}{\textbf{\checkmark}} \textcolor{greencheck}{#1}}
\newcommand{\highlightred}[1]{\textcolor{redcorrect}{\textbf{#1}}}
\newcommand{\highlightblue}[1]{\textcolor{bluecorrect}{\textbf{#1}}}
\newcommand{\tipbox}[1]{\fcolorbox{bluecorrect}{softyellow}{\parbox{0.88\textwidth}{\centering #1}}}
\newcommand{\steptitle}[1]{\textcolor{steppurple}{\textbf{#1}}}
\newcommand{\keypoint}[1]{\textcolor{deepblue}{\textbf{#1}}}

\newcommand{\stepbox}[1]{%
	\begin{center}
		\fcolorbox{darkgray}{white}{%
			\parbox{0.90\textwidth}{\vspace{0.1cm} #1 \vspace{0.1cm}}%
		}
	\end{center}
}

\begin{document}

% ================================================================
% SLIDE 1: TITLE
% ================================================================
\begin{frame}
\centering
\vspace{0.8cm}
{\fontsize{20}{24}\selectfont \textbf{Money Laundering and Financial Crime}}\\[0.2cm]
{\fontsize{15}{18}\selectfont Complete Tutorial Guide}\\[0.4cm]
{\fontsize{12}{14}\selectfont Step-by-Step Learning for Every Topic}\\[0.4cm]
{\fontsize{10}{12}\selectfont Compliance Training Department}
\vspace{0.3cm}
\rule{5cm}{1pt}
\end{frame}

% ================================================================
% SLIDE 2: TUTORIAL 1 — GLOBAL CHALLENGE
% ================================================================
\begin{frame}
\frametitle{\fontsize{14}{17}\selectfont TUTORIAL 1: Understanding the Global Challenge}

\small
\begin{block}{\fontsize{11}{13}\selectfont What We Need to Understand}
\footnotesize
Fighting financial crime globally is extremely difficult. Why? Because \textbf{different countries have different laws}. What is a crime in one country might be completely legal in another.
\end{block}

\vspace{0.03cm}

\stepbox{\footnotesize
\steptitle{Step 1:} Recognize the problem of legal divergence. Every country has its own legal system, its own definitions of crimes, and its own cultural values. This means there is no single global standard for what constitutes financial crime. \\
\steptitle{Step 2:} See how criminals exploit this. They \textbf{shift their activities} to jurisdictions with weaker laws or different definitions. They operate in the \textbf{gaps} between legal systems. \\
\steptitle{Step 3:} Understand the challenge of mutual legal assistance. When crime crosses borders, countries need to cooperate. But this requires complex \textbf{mutual legal assistance treaties (MLATs)}. These are slow, bureaucratic, and often ineffective. \\
\steptitle{Step 4:} Recognize that technology is not the enemy. Modern technology like AI, machine learning, and blockchain analytics \textbf{helps} detect crime. The problem is not obsolete technology — it is fragmented human systems. \\
\steptitle{Step 5:} Understand the consequence. Criminals can exploit legal gaps across borders. They can move money through countries with weak enforcement. This makes global AML efforts incredibly challenging.}

\vspace{0.08cm}
\tipbox{\footnotesize \textbf{CORE LESSON:} The problem is not technology — it is \textit{legal fragmentation}. Different countries, different rules. That is what makes global AML so hard.}
\end{frame}

% ================================================================
% SLIDE 3: TUTORIAL 2 — CULTURAL HOSPITALITY
% ================================================================
\begin{frame}
\frametitle{\fontsize{14}{17}\selectfont TUTORIAL 2: Cultural Hospitality and Bribery}

\small
\begin{block}{\fontsize{11}{13}\selectfont What We Need to Understand}
\footnotesize
In some cultures, expensive gifts and lavish entertainment are expected in business relationships. But where does hospitality end and bribery begin? This is a critical question for compliance.
\end{block}

\vspace{0.03cm}

\stepbox{\footnotesize
\steptitle{Step 1:} Recognize the cultural challenge. In many countries, offering gifts, meals, and entertainment to business partners and government officials is a \textbf{normal part of building relationships}. Employees may feel pressure to participate. \\
\steptitle{Step 2:} Understand the risk of blurring lines. Without clear guidance, employees do not know \textbf{where the line is}. They may offer excessive hospitality thinking it is "cultural" — when it is actually a \textbf{bribe}. \\
\steptitle{Step 3:} See the danger of inadvertent facilitation. The organization does not \textit{intend} to bribe. But the action still constitutes bribery under anti-corruption laws. This is called \textbf{inadvertent facilitation}. \\
\steptitle{Step 4:} Learn the solution. Organizations must set \textbf{clear monetary limits} on gifts and hospitality. They should require \textbf{pre-approval} for any entertainment of foreign officials. \\
\steptitle{Step 5:} Provide training. Employees need to understand the \textbf{distinction} between cultural expectations and corruption. They need to know what is acceptable and what is not.}

\vspace{0.08cm}
\tipbox{\footnotesize \textbf{CORE LESSON:} Never assume cultural expectations override compliance. Define the line clearly. The risk is \textit{inadvertent facilitation} — and it is dangerous.}
\end{frame}

% ================================================================
% SLIDE 4: TUTORIAL 3 — THE THREE STAGES
% ================================================================
\begin{frame}
\frametitle{\fontsize{14}{17}\selectfont TUTORIAL 3: The Three Stages of Money Laundering}

\small
\begin{block}{\fontsize{11}{13}\selectfont What We Need to Understand}
\footnotesize
Money laundering happens in three stages. Each stage serves a different purpose. Understanding where different techniques fit is essential for detection.
\end{block}

\vspace{0.03cm}

\stepbox{\footnotesize
\steptitle{Step 1:} Understand the first stage — \textbf{Placement}. This is where dirty cash enters the financial system. Methods include physical deposits, currency exchanges, and buying casino chips. The cash is still "dirty" and visible. \\
\steptitle{Step 2:} Understand the second stage — \textbf{Layering}. This is where the money is moved around to hide its origin. Methods include multiple transfers, complex structures, and \textbf{shell companies}. This is where the money gets \textbf{hidden}. \\
\steptitle{Step 3:} Understand the third stage — \textbf{Integration}. This is where the money is used as if it is "clean." Methods include investments, property purchases, and business funding. The money now appears legitimate. \\
\steptitle{Step 4:} See the role of shell companies. Shell companies are used in \textbf{layering}. They create \textbf{complex chains} of transactions. They \textbf{obscure beneficial ownership}. They make illegal money look like legitimate business income. \\
\steptitle{Step 5:} Remember the logic. \textbf{Placement} = putting money in. \textbf{Layering} = hiding where it came from. \textbf{Integration} = using it as clean money.}

\vspace{0.08cm}
\tipbox{\footnotesize \textbf{MEMORY TRICK:} \textit{"Placement = Put it in; Layering = Hide it; Integration = Use it."} Shell companies = \textbf{HIDING} = Layering.}
\end{frame}

% ================================================================
% SLIDE 5: TUTORIAL 4 — MOBILE ONBOARDING
% ================================================================
\begin{frame}
\frametitle{\fontsize{14}{17}\selectfont TUTORIAL 4: Mobile Onboarding and Identity Verification}

\small
\begin{block}{\fontsize{11}{13}\selectfont What We Need to Understand}
\footnotesize
Financial institutions face a challenge: they need strong identity verification for compliance, but customers want fast and easy mobile access. How do we balance these competing needs?
\end{block}

\vspace{0.03cm}

\stepbox{\footnotesize
\steptitle{Step 1:} Recognize the tension. There are \textbf{two competing needs}. Compliance requires \textit{strong} identity verification. Customers demand \textit{fast} and \textit{easy} access. These goals pull in opposite directions. \\
\steptitle{Step 2:} Understand the problem with choosing one. If you choose strong compliance, you create friction and lose customers. If you choose easy access, you increase fraud and regulatory risk. Neither option is acceptable. \\
\steptitle{Step 3:} Learn the solution — \textbf{biometric verification}. Fingerprint scanning and facial recognition provide \textbf{strong identity assurance} AND are \textbf{fast and easy} for customers. This is the best of both worlds. \\
\steptitle{Step 4:} Add a second layer — \textbf{digital onboarding review}. This ensures the \textit{process itself} is secure and compliant. It reviews the overall system, not just individual identities. \\
\steptitle{Step 5:} Achieve balance. This approach gives you \textbf{both}: strong identity verification \textbf{and} a smooth mobile experience. You do not have to choose.}

\vspace{0.08cm}
\tipbox{\footnotesize \textbf{CORE LESSON:} Do not choose between compliance and access — use \textit{technology} to achieve both. Biometrics are the answer.}
\end{frame}

% ================================================================
% SLIDE 6: TUTORIAL 5 — DEFENSIVE SARS
% ================================================================
\begin{frame}
\frametitle{\fontsize{14}{17}\selectfont TUTORIAL 5: Defensive Suspicious Activity Reports}

\small
\begin{block}{\fontsize{11}{13}\selectfont What We Need to Understand}
\footnotesize
A Suspicious Activity Report (SAR) is a critical tool for detecting financial crime. But when reports are filed for the wrong reasons, the system breaks down.
\end{block}

\vspace{0.03cm}

\stepbox{\footnotesize
\steptitle{Step 1:} Understand what a defensive SAR is. It is a report filed \textbf{not because there is genuine suspicion}, but because the filer wants to \textbf{avoid blame} later. It is a "cover your back" filing. \\
\steptitle{Step 2:} Recognize the problem. Defensive SARs do \textbf{not} help the system. They create \textbf{more work}, not less. They add to the compliance burden without adding value. \\
\steptitle{Step 3:} Understand the larger issue. Regulators receive \textbf{thousands} of SARs every day. Defensive filings \textbf{flood} the system with \textit{irrelevant} reports. \\
\steptitle{Step 4:} See the consequence. This creates \textbf{noise}. Real suspicious activity gets lost in the flood. It becomes harder for regulators to find \textbf{genuine} threats. \\
\steptitle{Step 5:} Apply the correct approach. File SARs \textbf{only when you have reasonable suspicion}. Do not file them to "cover your back." Defensive filings hurt the entire AML system.}

\vspace{0.08cm}
\tipbox{\footnotesize \textbf{CORE LESSON:} File SARs only when you have \textit{reasonable suspicion} — not to "cover yourself." Defensive filings hurt everyone.}
\end{frame}

% ================================================================
% SLIDE 7: TUTORIAL 6 — LAW ENFORCEMENT CALLS
% ================================================================
\begin{frame}
\frametitle{\fontsize{14}{17}\selectfont TUTORIAL 6: Responding to Law Enforcement}

\small
\begin{block}{\fontsize{11}{13}\selectfont What We Need to Understand}
\footnotesize
Financial crime investigators sometimes receive calls from law enforcement. These calls may be informal and lack documentation. How should investigators respond?
\end{block}

\vspace{0.03cm}

\stepbox{\footnotesize
\steptitle{Step 1:} Understand the scenario. A law enforcement officer calls, asks for sensitive information, but provides \textbf{no written documentation}. The call is informal. \\
\steptitle{Step 2:} Recognize the risks. Without documentation, there is \textbf{no official record}. The investigator could \textbf{misremember} details. The call could be \textbf{impersonation} — a fraudster pretending to be law enforcement. \\
\steptitle{Step 3:} Take immediate action. \textbf{Document the key points} of the conversation \textit{right away}. Write down who called, when, what was asked, and what was discussed. \\
\steptitle{Step 4:} Request formal confirmation. Ask the officer to send \textbf{written documentation} of the request. This provides verification and creates an official record. \\
\steptitle{Step 5:} Understand why this matters. This creates an \textbf{audit trail}. It protects both the investigator \textbf{and} the institution. If there is ever a dispute, there is documentation.}

\vspace{0.08cm}
\tipbox{\footnotesize \textbf{CORE LESSON:} If it is not documented, it did not happen. \textit{Always} get it in writing. Documentation protects everyone.}
\end{frame}

% ================================================================
% SLIDE 8: TUTORIAL 7 — CONCENTRATION ACCOUNTS
% ================================================================
\begin{frame}
\frametitle{\fontsize{14}{17}\selectfont TUTORIAL 7: Concentration Accounts and Oversight}

\small
\begin{block}{\fontsize{11}{13}\selectfont What We Need to Understand}
\footnotesize
Concentration accounts are internal accounts used to aggregate transactions. Without proper oversight, they become a vulnerability for money laundering.
\end{block}

\vspace{0.03cm}

\stepbox{\footnotesize
\steptitle{Step 1:} Understand what a concentration account is. It is an internal account used to \textbf{aggregate} multiple transactions before they are posted to individual customer accounts. Think of it as a "holding" account. \\
\steptitle{Step 2:} Recognize the concern. If a concentration account lacks \textbf{sufficient oversight}, it becomes a \textbf{black box}. No one knows exactly what is happening inside it. \\
\steptitle{Step 3:} See the vulnerability. Funds go \textit{in} and come \textit{out}. But it is \textbf{hard to trace} where individual funds went. You cannot tell which money belongs to which customer. \\
\steptitle{Step 4:} Understand the criminal opportunity. Criminals can use this \textbf{poor traceability} to move money undetected. They can "mix" illicit funds with legitimate funds. \\
\steptitle{Step 5:} Recognize the regulatory concern. The institution might \textbf{inadvertently facilitate money laundering}. Not intentionally — but because the account lacks proper controls. This is why regulators \textbf{scrutinize} concentration accounts closely.}

\vspace{0.08cm}
\tipbox{\footnotesize \textbf{CORE LESSON:} If you cannot \textit{trace} it, you cannot \textit{control} it. Oversight of concentration accounts is critical.}
\end{frame}

% ================================================================
% SLIDE 9: TUTORIAL 8 — PRIORITIZING REVIEWS
% ================================================================
\begin{frame}
\frametitle{\fontsize{14}{17}\selectfont TUTORIAL 8: Prioritizing Transaction Reviews}

\small
\begin{block}{\fontsize{11}{13}\selectfont What We Need to Understand}
\footnotesize
Not all transactions are equally risky. Compliance officers need to prioritize their reviews based on risk indicators. Understanding which indicators matter is essential.
\end{block}

\vspace{0.03cm}

\stepbox{\footnotesize
\steptitle{Step 1:} Understand the scenario. Two trades have passed basic monitoring. One is high-volume equities with normal volatility. The other is low-volume penny stocks with rapid pricing shifts. Which needs more review? \\
\steptitle{Step 2:} Recognize the risk indicators. \textbf{Low volume} and \textbf{high volatility} are risk indicators. They suggest that a trade is easier to manipulate. \\
\steptitle{Step 3:} Understand why penny stocks are risky. Penny stocks have \textbf{low trading volume} — a single trade can move the price significantly. They have \textbf{high volatility} — prices can swing wildly. \\
\steptitle{Step 4:} See the manipulation risk. Criminals use penny stocks to \textbf{create artificial activity}. They can pump and dump, or create fake trading patterns. This is much harder with high-volume, stable equities. \\
\steptitle{Step 5:} Apply the lesson. Do not assume all activity is equal. \textbf{Lower volume} + \textbf{higher volatility} = \textbf{higher risk}. Prioritize these for review.}

\vspace{0.08cm}
\tipbox{\footnotesize \textbf{CORE LESSON:} Do not assume all activity is equal. \textit{Lower volume} + \textit{higher volatility} = \textbf{higher risk}. Prioritize accordingly.}
\end{frame}

% ================================================================
% SLIDE 10: TUTORIAL 9 — FACIAL RECOGNITION
% ================================================================
\begin{frame}
\frametitle{\fontsize{14}{17}\selectfont TUTORIAL 9: Biometric Risks}

\small
\begin{block}{\fontsize{11}{13}\selectfont What We Need to Understand}
\footnotesize
Biometric authentication like facial recognition is powerful, but it comes with a unique risk. Understanding this risk is essential for security planning.
\end{block}

\vspace{0.03cm}

\stepbox{\footnotesize
\steptitle{Step 1:} Understand the comparison. \textbf{Passwords} can be changed if compromised. If someone steals your password, you simply reset it. \textbf{Biometrics} cannot be changed easily. You have \textbf{one face}. \\
\steptitle{Step 2:} Recognize the fundamental vulnerability. If your \textbf{biometric data} (face scan, fingerprint) is compromised, you \textbf{cannot reset it}. \\
\steptitle{Step 3:} Understand the permanence. You have \textit{one} face. \textit{One} set of fingerprints. Once stolen, that data is \textbf{permanently} exposed. You cannot change your face like you change a password. \\
\steptitle{Step 4:} See the consequence. This is the \textbf{fundamental vulnerability} of biometric authentication. It is powerful and convenient, but if it is hacked, the damage is \textbf{permanent}. \\
\steptitle{Step 5:} Apply the lesson. Biometric data must be \textbf{protected even more carefully} than passwords. Because if it is stolen, there is no way to "reset" it.}

\vspace{0.08cm}
\tipbox{\footnotesize \textbf{CORE LESSON:} Passwords are replaceable. Biometrics are \textit{not}. That is why biometric data must be \textbf{protected even more carefully}.}
\end{frame}

% ================================================================
% SLIDE 11: TUTORIAL 10 — COMPLIANCE CULTURE
% ================================================================
\begin{frame}
\frametitle{\fontsize{14}{17}\selectfont TUTORIAL 10: Building a Compliance Culture}

\small
\begin{block}{\fontsize{11}{13}\selectfont What We Need to Understand}
\footnotesize
Having policies on paper is not enough. Organizations need to build a culture where employees genuinely believe in compliance. This is the foundation of effective AML.
\end{block}

\vspace{0.03cm}

\stepbox{\footnotesize
\steptitle{Step 1:} Recognize the problem. Many organizations have \textbf{robust policies} on paper. But employees do not \textbf{engage} with compliance. They see it as a \textbf{barrier} to their work. \\
\steptitle{Step 2:} Understand the gap. Having \textbf{policies on paper} is not enough. Compliance must be \textbf{embedded in the culture}. Employees must \textbf{believe} in compliance, not just follow rules. \\
\steptitle{Step 3:} See the consequence. When compliance is seen as a \textbf{barrier}, people will \textbf{find ways around it}. They will cut corners, ignore procedures, and rationalize their behavior. The policies fail. \\
\steptitle{Step 4:} Learn the solution. Organizations must \textbf{build a culture of compliance}. This means leadership must demonstrate commitment. Employees must be trained and motivated. Compliance must be seen as a \textbf{value}, not a burden. \\
\steptitle{Step 5:} Understand the truth. Compliance is not a set of rules — it is a \textit{culture}. If you do not build the culture, the rules will not work.}

\vspace{0.08cm}
\tipbox{\footnotesize \textbf{CORE LESSON:} Compliance is not a set of rules — it is a \textit{culture}. If you do not build the culture, the rules will not work.}
\end{frame}

% ================================================================
% SLIDE 12: TUTORIAL 11 — BOARD OVERSIGHT
% ================================================================
\begin{frame}
\frametitle{\fontsize{14}{17}\selectfont TUTORIAL 11: The Board's Role in AML}

\small
\begin{block}{\fontsize{11}{13}\selectfont What We Need to Understand}
\footnotesize
Board members have a specific role in AML oversight. They are not operators — they are overseers. Understanding this distinction is critical for effective governance.
\end{block}

\vspace{0.03cm}

\stepbox{\footnotesize
\steptitle{Step 1:} Understand the board's role. Board members are \textbf{oversight}, not \textit{operators}. They do not do the day-to-day work of compliance. They oversee the people who do. \\
\steptitle{Step 2:} Learn what effective oversight looks like. Board members should \textbf{receive regular reports} on financial crime risks and control effectiveness. They should \textbf{review whether senior management} is addressing audit findings. They should \textbf{challenge management} when performance metrics show weakness. \\
\steptitle{Step 3:} Understand what is NOT the board's role. Board members should NOT \textbf{participate in SAR investigations}. This is operational work. Boards should not be involved in operational details. \\
\steptitle{Step 4:} See the importance of challenging. Board members should ask \textbf{tough questions}. Is management acting on audit findings? Are risks increasing or decreasing? Are controls effective? \\
\steptitle{Step 5:} Apply the principle. Directors oversee, they do not operate. Their job is to \textit{challenge} and \textit{question} — not to \textit{do}.}

\vspace{0.08cm}
\tipbox{\footnotesize \textbf{CORE LESSON:} Directors oversee, they do not operate. Their job is to \textit{challenge} and \textit{question} — not to \textit{do}.}
\end{frame}

% ================================================================
% SLIDE 13: TUTORIAL 12 — TERRORIST FINANCING
% ================================================================
\begin{frame}
\frametitle{\fontsize{14}{17}\selectfont TUTORIAL 12: Terrorist Financing Techniques}

\small
\begin{block}{\fontsize{11}{13}\selectfont What We Need to Understand}
\footnotesize
Terrorists use a variety of techniques to move money across borders. Some are high-tech, some are low-tech. Understanding all of them is essential for detection.
\end{block}

\vspace{0.03cm}

\stepbox{\footnotesize
\steptitle{Step 1:} Learn the first technique — \textbf{Informal value transfer systems}. Hawala and Hundi are trust-based systems with no paper trail. They are almost impossible to detect. \\
\steptitle{Step 2:} Learn the second technique — \textbf{Physical transportation of cash}. Terrorists often simply \textbf{carry cash} across borders. This is old-school but \textbf{still highly effective}. No electronic trail. Works in regions with weak border controls. \\
\steptitle{Step 3:} Learn the third technique — \textbf{Fake charities}. Terrorists establish charities that appear legitimate. They move money under the cover of humanitarian work. This provides a plausible explanation for fund transfers. \\
\steptitle{Step 4:} Understand why these techniques work. They exploit \textbf{gaps} in the financial system and border controls. They use \textbf{trust} instead of documentation. They hide behind \textbf{legitimate appearances}. \\
\steptitle{Step 5:} Remember the lesson. Terrorists use \textbf{both high-tech and no-tech methods}. Do not overlook the simple ones like carrying cash.}

\vspace{0.08cm}
\tipbox{\footnotesize \textbf{CORE LESSON:} Terrorists use both \textit{high-tech} and \textit{no-tech} methods. Do not overlook the simple ones like carrying cash.}
\end{frame}

% ================================================================
% SLIDE 14: TUTORIAL 13 — HUMAN TRAFFICKING
% ================================================================
\begin{frame}
\frametitle{\fontsize{14}{17}\selectfont TUTORIAL 13: Human Trafficking Laundering}

\small
\begin{block}{\fontsize{11}{13}\selectfont What We Need to Understand}
\footnotesize
Human trafficking generates significant illicit proceeds. Traffickers use specific techniques to launder this money. Understanding these techniques is essential for detection.
\end{block}

\vspace{0.03cm}

\stepbox{\footnotesize
\steptitle{Step 1:} Learn the first technique — \textbf{Mobile payment and remittance systems}. Traffickers use mobile money because it is \textbf{hard to trace}, \textbf{fast}, and \textbf{cross-border}. They can move \textbf{small amounts frequently} without triggering alerts. \\
\steptitle{Step 2:} Learn the second technique — \textbf{Weak cross-border enforcement cooperation}. Traffickers operate across borders. Lack of cooperation between countries means they can evade detection. Jurisdictional gaps are exploited. \\
\steptitle{Step 3:} Understand why these are the key enablers. Mobile payments are \textbf{accessible} — anyone can use them. Weak cooperation means traffickers can \textbf{move freely} across borders without fear. \\
\steptitle{Step 4:} Understand what is NOT the primary enabler. Counterfeit documents are used in trafficking, but they are not the \textbf{primary laundering enabler}. The financial techniques are what matter for laundering. \\
\steptitle{Step 5:} Remember the lesson. Traffickers use \textit{accessible} systems (mobile money) and exploit \textit{jurisdictional gaps} (weak cooperation).}

\vspace{0.08cm}
\tipbox{\footnotesize \textbf{CORE LESSON:} Traffickers use \textit{accessible} systems and exploit \textit{jurisdictional gaps}. Mobile payments are a key enabler.}
\end{frame}

% ================================================================
% SLIDE 15: SUMMARY
% ================================================================
\begin{frame}
\frametitle{\fontsize{14}{17}\selectfont SUMMARY: Key Lessons from This Tutorial}

\small
\begin{block}{\fontsize{11}{13}\selectfont What We Have Learned:}
\footnotesize
\begin{enumerate}
    \item \textbf{Global Challenge:} Legal divergence is the real problem — not technology. Different countries, different rules.
    \item \textbf{Cultural Hospitality:} Cultural expectations can blur the line between hospitality and bribery. Define the line clearly.
    \item \textbf{Three Stages:} Shell companies belong in \textbf{layering}, not placement. Layering = hiding.
    \item \textbf{Mobile Onboarding:} Use biometrics to balance strong compliance with smooth customer access.
    \item \textbf{Defensive SARs:} They flood the system with noise. File only when you have reasonable suspicion.
    \item \textbf{Law Enforcement Calls:} Document everything immediately. Always get written confirmation.
    \item \textbf{Concentration Accounts:} If you cannot trace it, you cannot control it. Oversight is critical.
    \item \textbf{Prioritizing Reviews:} Low volume + high volatility = high risk. Prioritize accordingly.
    \item \textbf{Facial Recognition:} Biometrics cannot be reset like passwords. Protect them more carefully.
    \item \textbf{Compliance Culture:} Policies are not enough. Build a culture of compliance.
    \item \textbf{Board Oversight:} Directors oversee — they do not operate. Challenge and question.
    \item \textbf{Terrorist Financing:} Do not overlook simple methods like carrying cash.
    \item \textbf{Human Trafficking:} Mobile payments are a key laundering enabler.
\end{enumerate}
\end{block}

\vspace{0.05cm}
\tipbox{\footnotesize \textbf{FINAL LESSON:} Every topic we studied has one thing in common — \textit{understanding the underlying principle} is more important than memorizing facts. Know \textit{why} and you will know \textit{what} to do.}
\end{frame}

% ================================================================
% SLIDE 16: THANK YOU
% ================================================================
\begin{frame}
\centering
\vspace{2.5cm}
{\fontsize{22}{26}\selectfont \textbf{Thank You}} \\[0.3cm]
{\fontsize{14}{17}\selectfont Keep learning. Keep growing.} \\[0.2cm]
\rule{3cm}{1pt} \\[0.2cm]
{\fontsize{10}{12}\selectfont \textit{Compliance is not about perfection — it is about continuous improvement.}}
\vspace{0.5cm}
{\fontsize{9}{11}\selectfont \textcolor{darkgray}{End of Tutorial}}
\end{frame}

\end{document}